% Figures: cayley_fig.tex (the action on the cube and the cycle), placement.tex (four trees).
% Needs in the preamble: \newtheorem{theorem}{Theorem} and \newtheorem{corollary}{Corollary}.

% cayley_fig.tex: a template with ranks 0, 1 and 2 moved by the placement's group action on four Cayley graphs:
% (a) F_2^3 translated by 101, (b) Z_8 with every bank rotated by 4, (c) F_2^4 translated by 1100, (d) Z_16 with every bank rotated by 8. Edges are
% coloured by the generator they apply. Self-contained; generated by make_family.py and make_tesseract_octagon.py.
% Needs: \usepackage{tikz,amssymb} \usetikzlibrary{arrows.meta, backgrounds}
\begin{figure*}[!t]
\centering
\definecolor{gA}{HTML}{E69F00}\definecolor{gB}{HTML}{56B4E9}\definecolor{gC}{HTML}{009E73}\definecolor{gD}{HTML}{CC79A7}%
\tikzset{
  vx/.style={circle, draw=black!80, line width=0.6pt, minimum size=2.5mm, inner sep=0pt},
  ot/.style={vx, fill=black!10}, tm/.style={ot}, im/.style={ot},
  lb/.style={font=\fontsize{5.5}{6}\selectfont, text=black!80, inner sep=0.5pt},
  g0/.style={draw=gA, line width=0.75pt}, g1/.style={draw=gB, line width=0.75pt},
  g2/.style={draw=gC, line width=0.75pt}, g3/.style={draw=gD, line width=0.75pt},
  ac/.style={draw=black, line width=0.55pt, -{Stealth[length=1.6mm, width=1.25mm]}, shorten <=-0.4pt, shorten >=-0.9pt}}%
\begin{tabular}{@{}c@{\hspace{5mm}}c@{\hspace{5mm}}c@{\hspace{5mm}}c@{}}
\begin{tikzpicture}[baseline=(current bounding box.center)]
  \node[tm] (c0) at (1.294,-2.156) {};
  \node[tm] (c1) at (2.731,-2.156) {};
  \node[tm] (c2) at (1.294,-0.862) {};
  \node[ot] (c3) at (2.731,-0.862) {};
  \node[im] (c4) at (1.869,-1.725) {};
  \node[im] (c5) at (3.306,-1.725) {};
  \node[ot] (c6) at (1.869,-0.431) {};
  \node[im] (c7) at (3.306,-0.431) {};
  \begin{scope}[on background layer]
  \draw[g0] (c0) -- (c1);
  \draw[g1] (c0) -- (c2);
  \draw[g2] (c0) -- (c4);
  \draw[g1] (c1) -- (c3);
  \draw[g2] (c1) -- (c5);
  \draw[g0] (c2) -- (c3);
  \draw[g2] (c2) -- (c6);
  \draw[g2] (c3) -- (c7);
  \draw[g0] (c4) -- (c5);
  \draw[g1] (c4) -- (c6);
  \draw[g1] (c5) -- (c7);
  \draw[g0] (c6) -- (c7);
  \draw[ac] (c2) .. controls (1.495,0.345) and (3.105,0.230) .. (c7);
  \draw[ac] (c0) .. controls (1.495,-2.817) and (3.105,-2.875) .. (c5);
  \draw[ac] (c1) .. controls (2.645,-2.932) and (2.012,-2.932) .. (c4);
  \end{scope}
  \node[lb] at (1.043,-2.332) {000};
  \node[lb] at (2.346,-2.317) {001};
  \node[lb] at (0.978,-0.756) {010};
  \node[lb] at (2.616,-0.613) {011};
  \node[lb] at (1.551,-1.622) {100};
  \node[lb] at (3.622,-1.831) {101};
  \node[lb] at (2.080,-0.227) {110};
  \node[lb] at (3.557,-0.256) {111};
\end{tikzpicture}
&
\begin{tikzpicture}[baseline=(current bounding box.center)]
  \node[tm] (r0) at (0.000,1.050) {};
  \node[tm] (r1) at (0.742,0.742) {};
  \node[tm] (r2) at (1.050,0.000) {};
  \node[tm] (r3) at (0.742,-0.742) {};
  \node[tm] (r4) at (0.000,-1.050) {};
  \node[tm] (r5) at (-0.742,-0.742) {};
  \node[tm] (r6) at (-1.050,-0.000) {};
  \node[tm] (r7) at (-0.742,0.742) {};
  \begin{scope}[on background layer]
  \draw[g0] (0,0) circle (1.05cm);
  \draw[ac] (r0) to[bend right=26] (r4);
  \draw[ac] (r1) to[bend right=26] (r5);
  \draw[ac] (r2) to[bend right=26] (r6);
  \draw[ac] (r3) to[bend right=26] (r7);
  \draw[ac] (r4) to[bend right=26] (r0);
  \draw[ac] (r5) to[bend right=26] (r1);
  \draw[ac] (r6) to[bend right=26] (r2);
  \draw[ac] (r7) to[bend right=26] (r3);
  \end{scope}
  \node[lb] at (0.000,1.320) {0};
  \node[lb] at (0.951,0.933) {1};
  \node[lb] at (1.355,0.000) {2};
  \node[lb] at (0.951,-0.933) {3};
  \node[lb] at (0.000,-1.320) {4};
  \node[lb] at (-0.951,-0.933) {5};
  \node[lb] at (-1.355,-0.000) {6};
  \node[lb] at (-0.951,0.933) {7};
\end{tikzpicture}
&
\begin{tikzpicture}[baseline=(current bounding box.center)]
  \node[tm] (n0) at (0.619,-0.619) {};
  \node[tm] (n1) at (2.100,0.000) {};
  \node[tm] (n2) at (-0.875,-0.000) {};
  \node[ot] (n3) at (0.619,0.619) {};
  \node[ot] (n4) at (0.000,0.875) {};
  \node[ot] (n5) at (1.485,1.485) {};
  \node[ot] (n6) at (-1.485,1.485) {};
  \node[ot] (n7) at (0.000,2.100) {};
  \node[ot] (n8) at (0.000,-2.100) {};
  \node[ot] (n9) at (1.485,-1.485) {};
  \node[ot] (n10) at (-1.485,-1.485) {};
  \node[ot] (n11) at (0.000,-0.875) {};
  \node[im] (n12) at (-0.619,-0.619) {};
  \node[im] (n13) at (0.875,0.000) {};
  \node[im] (n14) at (-2.100,-0.000) {};
  \node[ot] (n15) at (-0.619,0.619) {};
  \begin{scope}[on background layer]
  \draw[g3] (n0) -- (n8);
  \draw[g2] (n0) -- (n4);
  \draw[g1] (n0) -- (n2);
  \draw[g0] (n0) -- (n1);
  \draw[g3] (n1) -- (n9);
  \draw[g2] (n1) -- (n5);
  \draw[g1] (n1) -- (n3);
  \draw[g3] (n2) -- (n10);
  \draw[g2] (n2) -- (n6);
  \draw[g0] (n2) -- (n3);
  \draw[g3] (n3) -- (n11);
  \draw[g2] (n3) -- (n7);
  \draw[g3] (n4) -- (n12);
  \draw[g1] (n4) -- (n6);
  \draw[g0] (n4) -- (n5);
  \draw[g3] (n5) -- (n13);
  \draw[g1] (n5) -- (n7);
  \draw[g3] (n6) -- (n14);
  \draw[g0] (n6) -- (n7);
  \draw[g3] (n7) -- (n15);
  \draw[g2] (n8) -- (n12);
  \draw[g1] (n8) -- (n10);
  \draw[g0] (n8) -- (n9);
  \draw[g2] (n9) -- (n13);
  \draw[g1] (n9) -- (n11);
  \draw[g2] (n10) -- (n14);
  \draw[g0] (n10) -- (n11);
  \draw[g2] (n11) -- (n15);
  \draw[g1] (n12) -- (n14);
  \draw[g0] (n12) -- (n13);
  \draw[g1] (n13) -- (n15);
  \draw[g0] (n14) -- (n15);
  \draw[ac] (n0) to[bend right=32] (n12);
  \draw[ac] (n1) to[bend right=32] (n13);
  \draw[ac] (n2) to[bend right=32] (n14);
  \end{scope}
  \node[lb] at (0.861,-0.810) {0000};
  \node[lb] at (2.472,0.000) {0001};
  \node[lb] at (-1.212,-0.103) {0010};
  \node[lb] at (0.861,0.810) {0011};
  \node[lb] at (0.000,1.145) {0100};
  \node[lb] at (1.727,1.676) {0101};
  \node[lb] at (-1.727,1.676) {0110};
  \node[lb] at (0.000,2.370) {0111};
  \node[lb] at (0.000,-2.370) {1000};
  \node[lb] at (1.727,-1.676) {1001};
  \node[lb] at (-1.727,-1.676) {1010};
  \node[lb] at (0.000,-1.145) {1011};
  \node[lb] at (-0.861,-0.810) {1100};
  \node[lb] at (1.212,-0.103) {1101};
  \node[lb] at (-2.472,-0.000) {1110};
  \node[lb] at (-0.861,0.810) {1111};
\end{tikzpicture}
&
\begin{tikzpicture}[baseline=(current bounding box.center)]
  \node[tm] (s0) at (0.000,1.350) {};
  \node[tm] (s1) at (0.517,1.247) {};
  \node[tm] (s2) at (0.955,0.955) {};
  \node[tm] (s3) at (1.247,0.517) {};
  \node[tm] (s4) at (1.350,0.000) {};
  \node[tm] (s5) at (1.247,-0.517) {};
  \node[tm] (s6) at (0.955,-0.955) {};
  \node[tm] (s7) at (0.517,-1.247) {};
  \node[tm] (s8) at (0.000,-1.350) {};
  \node[tm] (s9) at (-0.517,-1.247) {};
  \node[tm] (s10) at (-0.955,-0.955) {};
  \node[tm] (s11) at (-1.247,-0.517) {};
  \node[tm] (s12) at (-1.350,-0.000) {};
  \node[tm] (s13) at (-1.247,0.517) {};
  \node[tm] (s14) at (-0.955,0.955) {};
  \node[tm] (s15) at (-0.517,1.247) {};
  \begin{scope}[on background layer]
  \draw[g0] (0,0) circle (1.35cm);
  \draw[ac] (s0) to[bend right=24] (s8);
  \draw[ac] (s1) to[bend right=24] (s9);
  \draw[ac] (s2) to[bend right=24] (s10);
  \draw[ac] (s3) to[bend right=24] (s11);
  \draw[ac] (s4) to[bend right=24] (s12);
  \draw[ac] (s5) to[bend right=24] (s13);
  \draw[ac] (s6) to[bend right=24] (s14);
  \draw[ac] (s7) to[bend right=24] (s15);
  \draw[ac] (s8) to[bend right=24] (s0);
  \draw[ac] (s9) to[bend right=24] (s1);
  \draw[ac] (s10) to[bend right=24] (s2);
  \draw[ac] (s11) to[bend right=24] (s3);
  \draw[ac] (s12) to[bend right=24] (s4);
  \draw[ac] (s13) to[bend right=24] (s5);
  \draw[ac] (s14) to[bend right=24] (s6);
  \draw[ac] (s15) to[bend right=24] (s7);
  \end{scope}
  \node[lb] at (0.000,1.620) {0};
  \node[lb] at (0.625,1.497) {1};
  \node[lb] at (1.163,1.146) {2};
  \node[lb] at (1.526,0.620) {3};
  \node[lb] at (1.655,0.000) {4};
  \node[lb] at (1.526,-0.620) {5};
  \node[lb] at (1.163,-1.146) {6};
  \node[lb] at (0.625,-1.497) {7};
  \node[lb] at (0.000,-1.620) {8};
  \node[lb] at (-0.625,-1.497) {9};
  \node[lb] at (-1.174,-1.146) {10};
  \node[lb] at (-1.545,-0.620) {11};
  \node[lb] at (-1.677,-0.000) {12};
  \node[lb] at (-1.545,0.620) {13};
  \node[lb] at (-1.174,1.146) {14};
  \node[lb] at (-0.628,1.497) {15};
\end{tikzpicture}
\\[2mm]
{\footnotesize (a) $\mathbb{F}_2^{\,3}$, rotated by $101$} &
{\footnotesize (b) $\mathbb{Z}_8$, every bank rotated by $4$} &
{\footnotesize (c) $\mathbb{F}_2^{\,4}$, rotated by $1100$} &
{\footnotesize (d) $\mathbb{Z}_{16}$, every bank rotated by $8$}
\end{tabular}
\caption{The action of one element on the Cayley graphs of four groups of banks: on a template with ranks 0, 1 and 2 in (a) and (c), and on every bank in (b) and (d).
The illustrations are inspired by \cite{macauley2024revealinghiddenbeautyfinite}.}
\label{fig:cayley}
\end{figure*}

\section{Theory}\label{sec:theory}

\paragraph{\underline{Setup}}
A request carries a $W$-bit key $x$, which we split into its \emph{anchor} $hi$, the top
$W - K$ bits, and its \emph{body} $lo$, the low $K$ bits, so that $x = 2^K \cdot hi + lo$. Our split fixes two limits of Das and Pinotti~\cite{das2000mappings}, whose colours a tree
of height $h$ whole: its $h + 1$ modules grow with the tree, and its balance holds only
within a constant factor.

The $B = 2^b$ banks are the elements of an abelian group $(G, \circ)$ of order $B$, with
Cayley graph $\Gamma$. A template reads at most $K$ banked nodes, one per level of the
body, so we take $B \ge K$ throughout. Each $c \in G$ acts on $G$ by
$\phi_c(r) = c \circ r$. This action is
\emph{regular}:
\begin{enumerate}
	\item[\textbf{(P1)}] $\forall c \in G:\ \phi_c \in \mathrm{Sym}(G)$, with $\phi_c^{-1} = \phi_{c^{-1}}$;
	\item[\textbf{(P2)}] $\forall r, t \in G\ \ \exists!\, c \in G:\ \phi_c(r) = t$, namely $c = t \circ r^{-1}$.
\end{enumerate}
Two groups have direct hardware implementations, drawn as Cayley graphs in
Fig.~\ref{fig:cayley}: (i)~$(\mathbb{F}_2^{\,b}, \oplus)$, with $\phi_c(r) = c \oplus r$, a
hypercube $Q_b$ realised as $b$ XOR gates; and (ii)~$(\mathbb{Z}_B, +)$, with
$\phi_c(r) = c + r \bmod B$, a cycle $C_B$ realised as a $b$-bit adder.

\paragraph{\underline{Placement}}
A placement is given by two maps: a \emph{rank} $\rho$ from the nonzero bodies to $G$, and
a \emph{rotation} $g$ from the anchors to $G$. A node with $lo \ne 0$ is stored at address
$lo$ in bank
$$
	\mathrm{bank}(x) = \phi_{g(hi)}\bigl(\rho(lo)\bigr) = g(hi) \circ \rho(lo),
$$
and a node with $lo = 0$ is held in a register. We require \textbf{(P3)}~$\rho$ injective on every
template, and \textbf{(P4)}~$g$ a bijection, which needs exactly $B$ anchors, $W - K = b$. The rank
is whatever changes at every step of a template: the level $v_2(lo)$ for a Fenwick tree,
the depth for a heap or a trie.

\begin{theorem}[Conflict freedom]\label{thm:conflict}
	If $\rho$ is injective on every template, the banked nodes of every template lie in
	distinct banks.
\end{theorem}

\begin{proof}
	The banked nodes of a template share one anchor, so all are placed by the same bijection
	$\phi_{g(hi)}$, and their ranks are distinct.
\end{proof}

\paragraph{\underline{Balance}}
Let $\mu, \nu : G \to [0,1]$ be the distributions of $g(hi)$ and $\rho(lo)$ over the
banked nodes; they may model workload accesses or the storage of sparse and compacted
structures. Assuming independence,\footnote{This fails for coupled datasets, treated in later sections.} bank $t$ receives the fraction
$p(t) = (\mu \ast \nu)(t) = \sum_{c \in G} \mu(c)\,\nu(c^{-1} \circ t)$ of banked nodes,
with RMS imbalance
$\varepsilon_p = \sqrt{\frac{1}{B}\sum_t \bigl(B\,p(t) - 1\bigr)^2}$. To make
analysis easier, we introduce the Fourier
transform $\hat\mu(k) = \sum_{t} \mu(t)\,\overline{\chi_k(t)}$, with $\chi_k$ ranging over
the characters of $G$, turning this convolution into the product
$\hat p(k) = \hat\mu(k)\,\hat\nu(k)$. On $\mathbb{F}_2^{\,b}$ it is the Walsh--Hadamard
transform, $\chi_k(t) = (-1)^{k \cdot t}$; on $\mathbb{Z}_B$ it is the DFT,
$\chi_k(t) = e^{2\pi i k t/B}$.

\begin{theorem}[Balance]\label{thm:balance}
	If $g(hi)$ and $\rho(lo)$ are independent, then
	$\varepsilon_p = \sqrt{\sum_{k \neq 0} |\hat\mu(k)|^2\,|\hat\nu(k)|^2}$.
\end{theorem}

\begin{proof}
	Let $f(t) = B\,p(t) - 1$. Then
	$$
		\begin{aligned}
			\varepsilon_p^2
			 & = \frac{1}{B}\sum_t |f(t)|^2
			\overset{(a)}{=} \frac{1}{B^2}\sum_k |\hat f(k)|^2 \\
			 & \overset{(b)}{=} \sum_{k \neq 0} |\hat p(k)|^2
			\overset{(c)}{=} \sum_{k \neq 0} |\hat\mu(k)|^2\,|\hat\nu(k)|^2 ,
		\end{aligned}
	$$
	where (a) is Parseval's identity; (b) holds because, by linearity,
	$\hat f(k) = B\,\hat p(k) - \hat 1(k)$, and $\hat 1(k) = \sum_t \overline{\chi_k(t)}$ is
	$B$ at $k = 0$ and $0$ otherwise, so $\hat f(0) = B(\hat p(0) - 1) = 0$; and (c) is the
	product rule.
\end{proof}

Imbalance thus survives at a frequency only if both $g(hi)$ and $\rho(lo)$ are uneven
there; making either one uniform removes it entirely.

\begin{corollary}[Exact balance]\label{cor:exact}
	If every anchor holds every body, every bank holds exactly $2^K - 1$ nodes.
\end{corollary}

\begin{proof}
	Every anchor--body pair occurs once, so $g(hi)$ and $\rho(lo)$ are independent and, by
	(P4), $g(hi)$ is uniform. Hence $\hat\mu(k) = 0$ for $k \neq 0$, so $\varepsilon_p = 0$
	and each bank receives $1/B$ of the $B(2^K - 1)$ banked nodes.
\end{proof}

\paragraph{\underline{Non-bijective rotations}}
A cache holds only some anchors, so its $g$ need not be one-to-one, and anchors with the
same rotation place each body in the same bank at the same address.

\begin{theorem}[Non-bijective rotation]\label{thm:cache}
	Let $g$ map the anchors onto $G$, and let an anchor whose rotation is in use overwrite the
	one held there. Every resident template is conflict-free, and if every anchor holds every
	body, then once every rotation is filled, every bank holds exactly $2^K - 1$ nodes.
\end{theorem}

\begin{proof}
	Two anchors with the same rotation store each body in the same bank at the same address,
	so at most one of them is resident, and $g$ is injective on the resident anchors.
	Theorem~\ref{thm:conflict} does not depend on $g$, so it holds for every resident
	template. Once every rotation is filled, $g$ is a bijection from the resident anchors onto
	$G$, and Corollary~\ref{cor:exact} applies to them.
\end{proof}

\paragraph{\underline{Linear placements}}
Theorem~\ref{thm:conflict} needs only $B \ge K$ banks while linear placements using XOR schemes in hardware, need exponentially more.

\begin{theorem}[Linear placements]\label{thm:linear}
	No linear placement $Ax \oplus u$ with fewer than $2^{K-1}$ banks is conflict-free on
	Fenwick trees.
\end{theorem}

\begin{proof}
	Within an anchor only the columns $a_0, \ldots, a_{K-1}$ of $A$ acting on the body matter.
	A query on a body with set bits $P$ visits $v_j$, the body with its bits below $j$
	cleared, for each $j \in P$, and for $j < k$ in $P$,
	$\mathrm{bank}(v_j) \oplus \mathrm{bank}(v_k) = \bigoplus_{i \in P,\, j \le i < k} a_i$.
	For any nonempty $T \subseteq \{0, \ldots, K - 2\}$, the body with set bits
	$T \cup \{K - 1\}$, with $j = \min T$ and $k = K - 1$, needs
	$\bigoplus_{i \in T} a_i \ne 0$. Hence $a_0, \ldots, a_{K-2}$ are linearly independent in
	$\mathbb{F}_2^{\,b}$, so $b \ge K - 1$ and $B \ge 2^{K-1}$.
\end{proof}
